\section{Three-Value Structure and Two-Value Degeneration of Interval Window Return Times}

Fix the window $E=(0,b)\subset\TT$ with $b\in(0,1)$. Endpoint conventions affect only sets of measure zero; below we work almost everywhere.

\subsection{Three-Value Partition Theorem}

\begin{theorem}[Three-Value Return Spectrum and Piecewise Constancy]\label{thm:three-value}
Take $i\ge 0$ such that
\[
\delta_i<b\le \delta_{i-1},
\]
and define
\[
K=\max\{k\ge 0:\ k\delta_i+\delta_{i+1}<b\}.
\]
Then the return-time function $r_E$ on $E$ takes at most three values, and there exists a partition into at most three interval pieces on which $r_E$ is constant. The set of possible values can be expressed as
\[
\mathrm{Spec}(E)\subseteq \left\{q_i,\ q_{i+1}-(K-1)q_i,\ q_{i+1}-Kq_i\right\},
\]
and $0\le K\le a_{i+1}-1$. Moreover, under suitable ordering, the values satisfy an additive closure relation: the middle value equals the sum of the two outer values.
\end{theorem}

\begin{proof}[Proof Outline]
See Appendix \ref{app:proof-thm31}. The proof uses return-tower decomposition: by comparing $b$ with the error hierarchy $(\delta_i,\delta_{i+1})$ we determine the possible tower heights; by rotation equidistance and window connectivity we prove each tower height corresponds to an interval return domain, yielding piecewise constancy and the three-value bound.
\end{proof}

\subsection{Necessary and Sufficient Conditions for Two-Value Degeneration}

\begin{theorem}[Two-Value Degeneration Criterion]\label{thm:two-value-criterion}
Under the notation of Theorem \ref{thm:three-value}, the following conditions are equivalent:
\begin{enumerate}
  \item $r_E$ takes exactly two values on $E$;
  \item The middle piece degenerates to zero length;
  \item $b$ satisfies the linear equation
  \[
  b=(K+1)\delta_i+\delta_{i+1};
  \]
  \item There exists an integer $c$ satisfying
  \[
  b=\delta_{i-1}-c\delta_i,\qquad 0\le c<a_{i+1}.
  \]
\end{enumerate}
\end{theorem}

\begin{proof}
By the three-piece structure of Theorem \ref{thm:three-value}, the middle piece degenerates when its endpoints coincide, which is equivalent to $b=(K+1)\delta_i+\delta_{i+1}$. Using $\delta_{i-1}=a_{i+1}\delta_i+\delta_{i+1}$ to rewrite gives
$b=\delta_{i-1}-c\delta_i$ where $c=a_{i+1}-(K+1)$, automatically satisfying $0\le c<a_{i+1}$. Proven.
\end{proof}

\subsection{Two-Value Degeneration and Induction Preserves Rotation}

\begin{theorem}[Dynamical Equivalence of Two-Value Degeneration]\label{thm:two-value-iet}
For an interval window $E=(0,b)$, the following propositions are equivalent:
\begin{enumerate}
  \item $r_E$ takes only two values on $E$;
  \item The induced map $T_{\theta,E}$ is a 2-IET;
  \item $(E,m_E,T_{\theta,E})$ is measure-conjugate to some irrational rotation system.
\end{enumerate}
\end{theorem}

\begin{proof}
(1)$\Rightarrow$(2): Two-value return time means $E$ decomposes into two pieces, each on which $T_{\theta,E}$ is a pure translation, hence 2-IET.\\
(2)$\Rightarrow$(3): A 2-IET is equivalent via endpoint identification and linear conjugacy to a circle rotation.\\
(3)$\Rightarrow$(1): The rotation of a two-piece section induces two-level return towers, hence only two values. Proven.
\end{proof}
